% ==============================================================================
% ПАРТИЯ B03-027: DISS-005-B1525 .. DISS-005-B1563
% Глава 3: Сбор осциллограмм, поиск и формирование единой БД
% Замечания: DISS-005-C0112 .. DISS-005-C0116 (5 шт.)
% Объекты: 9 шт. (4 макетные таблицы, 5 иллюстраций)
% ==============================================================================

% DISS-005-B1525
\section{Сбор и первичная обработка осциллограмм}\label{sec:ch3_data_collection}

% DISS-005-B1526
Начальным этапом формирования датасета является сбор исходных данных~-- файлов осциллограмм. Этот процесс требует не только доступа к разнообразным источникам, но и разработки инструментов для эффективного поиска, фильтрации и первичной систематизации собираемого материала.\par\vspace{0.2cm}

% DISS-005-B1527
\subsection{Источники данных и типы осциллограмм}\label{sec:ch3_data_sources}

% DISS-005-B1528
Осциллограммы для настоящего исследования были собраны из внутренних баз данных энергопредприятий, архивов пусконаладочных работ, а также из материалов, предоставленных для анализа аварийных ситуаций и оценки поведения устройств РЗА. На первом этапе сбора данных преобладали осциллограммы из электрических сетей классов напряжения 0,4--35~кВ с пусконалодочных работ. Как следствие, начальный датасет включает значительное количество так называемых «пустых» осциллограмм [95, 96], то есть записей, где отсутствуют ярко выраженные аварийные или переходные процессы, и фиксируется преимущественно шум, испытания с помощью специальных источников, таких как РЕТОМ, или нормальный режим работы (см. Рисунок~3.1). Несмотря на кажущуюся малоинформативность, такие осциллограммы важны для обучения моделей машинного обучения различению между нормальным режимом, помехами и действительными событиями.\par\vspace{0.2cm}

% DISS-005-TAB-072 (макетная таблица для рисунка 3.1)
\begin{figure}[htbp]
    \centering
    % DISS-005-B1529 / DISS-005-FIG-082
    \includegraphics[width=0.85\textwidth]{image85.jpeg}
    % DISS-005-B1530
    \caption[Пример «Пустой» осциллограммы]{Пример «Пустой» осциллограммы\commentnote{DISS-005-C0112}{Алексей Евдаков [2]}{Эта и многие другие картинки находятся в папке\newline\newline D:\textbackslash{}\textbackslash{}Учебное\textbackslash{}\textbackslash{}Аспирантура\textbackslash{}\textbackslash{}4) Научные исследования\textbackslash{}\textbackslash{}2) Публикация статей\textbackslash{}\textbackslash{}2024-2025\textbackslash{}\textbackslash{}(IEEE PCMP) Подготовка датасета\textbackslash{}\textbackslash{}Схемы и картинки\newline\newline Отсюда их можно взять и крректировать}.}\label{fig:ch3_empty_oscillogram}
\end{figure}

% DISS-005-B1531
Значительная часть осциллограмм в текущем датасете относится к работе систем быстродействующего автоматического ввода резерва (БАВР) на двухсекционных подстанциях (пример на Рисунок~3.2). Постепенно осуществляется расширение датасета за счет осциллограмм, отражающих работу традиционных устройств РЗА, а также данных из сетей более высоких классов напряжения (например, 110~кВ и выше).\par\vspace{0.2cm}

% DISS-005-B1532
Как уже отмечалось в разделе~3.2, одной из сложностей является разнообразие форматов осциллограмм. Большинство современных устройств РЗА записывают данные в различных версиях формата COMTRADE (файлы «.cfg + .dat» или единый «.cff»). Однако, помимо этого, в работе были использованы и другие форматы, такие как «.brs» (НПП Бреслер), «.dfr» (НПП ЭКРА) и др. Для обеспечения возможности дальнейшей унифицированной обработки была реализована поддержка чтения и первичной обработки основных из этих форматов.\par\vspace{0.2cm}

% DISS-005-B1533

% DISS-005-TAB-073, DISS-005-TAB-074 (макетные таблицы для рисунка 3.2)
\begin{figure}[htbp]
    \centering
    % DISS-005-B1534 / DISS-005-FIG-083
    \includegraphics[width=0.76\textwidth]{image86.jpeg}\par
    % DISS-005-B1535
    \vspace{2pt}{\small а) Секция 1}\par\vspace{8pt}
    % DISS-005-B1536 / DISS-005-FIG-084
    \includegraphics[width=0.76\textwidth]{image87.jpeg}\par
    % DISS-005-B1537
    \vspace{2pt}{\small б) Секция 2}\par\vspace{6pt}
    % DISS-005-B1538
    \caption[Пример осциллограммы срабатывания быстродействующего автоматического ввода резерва]{Пример осциллограммы срабатывания быстродействующего автоматического ввода резерва. Дискретный сигнал представляет используемую ручную разметку.}\label{fig:ch3_bavr_operation_example}
\end{figure}

% DISS-005-B1539
\subsection{Автоматизированный поиск и сбор осциллограмм}\label{sec:ch3_auto_search}

% DISS-005-B1540
Учитывая потенциально огромные объемы и распределенный характер хранения исходных данных (серверы предприятий, локальные компьютеры инженеров, архивы), ручной поиск и отбор осциллограмм являются неэффективными. Для решения этой и множества других задача на языке Python был разработан программный модуль (класс SearchOscillograms\commentnote{DISS-005-C0113}{Алексей Евдаков [2]}{Пока указываю функции, чтобы самому было проще, но в будущем, скорее всего будет удалено.}), реализующий автоматизированный поиск и копирование файлов осциллограмм с учетом их уникальности. Также для упрощения сбора и обработки данных иными лицами, была осуществлена модернизация частной программы APScilloscope предназначенной для просмотра осциллограмм в формате COMTRADE, а именно~-- добавление функции поиска осциллограммы (данная ПО предоставляется в свободном доступе компанией ООО «АПС» [XX\commentnote{DISS-005-C0114}{Алексей Евдаков [2]}{Ссылка на сайт, или что-то другое}\commentnote{DISS-005-C0115}{Алексей Евдаков [2]}{Подумать о том, стоит ли указывать это…}] занимающейся выпуском современного оборудования РЗА).\par\vspace{0.2cm}

% DISS-005-B1541
Логика работы данного модуля (copy\_new\_oscillograms) включает следующие ключевые шаги:\par\vspace{0.2cm}

% DISS-005-B1542
\textbf{Рекурсивный обход указанных директорий:} Программа осуществляет поиск файлов во всех вложенных папках, начиная с заданного корневого каталога. Это позволяет обрабатывать сложные структуры хранения данных.\par\vspace{0.2cm}

% DISS-005-B1543
\textbf{Поддержка различных форматов осциллограмм:} Функция идентифицирует файлы осциллограмм по их расширениям.\par\vspace{0.2cm}

% DISS-005-B1544
\textbf{Исключение дубликатов:} Для предотвращения многократного копирования одних и тех же осциллограмм, используется механизм хеширования. Вычисляется MD5-хеш содержимого файла данных (.dat), который затем используется как уникальный идентификатор осциллограммы. Перед копированием проверяется наличие данного хеша в словаре (copied\_hashes), содержащем информацию об уже обработанных файлах. Этот словарь может использоваться для инкрементального обновления коллекции осциллограмм, исключая повторное копирование уже найденных.\par\vspace{0.2cm}

% DISS-005-B1545
\textbf{Сохранение структуры директорий или копирование в единое хранилище:} Предусмотрена опция сохранения исходной структуры папок при копировании в целевую директорию (Preserve\_dir\_structure=True), что может быть полезно для сохранения контекста происхождения файла. Альтернативно, все найденные уникальные осциллограммы могут быть скопированы в одну общую папку, если пользователю важная конфиденциальность.\par\vspace{0.2cm}

% DISS-005-B1546
\textbf{Обработка архивных файлов:} Модуль способен распознавать и распаковывать распространенные форматы архивов (.zip, .7z, .rar) во временную директорию с последующим рекурсивным вызовом функции поиска уже внутри распакованного содержимого (\_process\_archive\_file). Это позволяет автоматически извлекать осциллограммы, хранящиеся в сжатом виде, что характерно для многих архивов данных.\par\vspace{0.2cm}

% DISS-005-B1547
\textbf{Ведение журнала обработки:} В процессе работы создаются и обновляются JSON-файлы (\_hash\_table.json и \_new\_hash\_table\_\{timestamp\}.json), содержащие информацию о хешах всех найденных уникальных осциллограмм и их исходных путях, а также список только что добавленных файлов. Это облегчает управление коллекцией и отслеживание ее пополнения.\par\vspace{0.2cm}

% DISS-005-B1548
\subsection{Формирование единой базы данных осциллограмм}\label{sec:ch3_single_db}

% DISS-005-B1549
В результате работы алгоритмов поиска и сбора формируется унифицированное хранилище осциллограмм. На данном этапе файлы могут быть структурированы по их исходному формату или по производителю оборудования, если такая информация доступна и релевантна. Важной частью этого этапа является создание и поддержка словаря уникальных файлов, который, как было упомянуто, ведется с использованием хеш-сумм.\par\vspace{0.2cm}

% DISS-005-B1550
Такая централизованная, хотя и первично обработанная, база данных становится отправной точкой для последующих, более глубоких этапов предобработки, таких как деперсонализация, унификация имен сигналов и нормализация (см. Рисунок~3.3) и в итоге будет сформирована открытый датасет (см. Рисунок~3.4).\par\vspace{0.2cm}

% DISS-005-TAB-075 (макетная таблица для рисунков 3.3 и 3.4)
\begin{figure}[htbp]
    \centering
    % DISS-005-B1551 / DISS-005-FIG-085
    \includegraphics[width=0.85\textwidth]{image88.png}
    % DISS-005-B1552
    \caption[Схема сбора и подготовки осциллограмм]{Схема сбора и подготовки осциллограмм.}\label{fig:ch3_data_collection_scheme}
\end{figure}

\begin{figure}[htbp]
    \centering
    % DISS-005-B1553 / DISS-005-FIG-086
    \includegraphics[width=0.6\textwidth]{image89.png}
    % DISS-005-B1554
    \caption[Структура публичного датасета]{Структура публичного датасета\commentnote{DISS-005-C0116}{Алексей Евдаков [2]}{Оно может меняться. Да и в целом вопрос, насколько это нужно здесь.}.}\label{fig:ch3_public_dataset_structure}
\end{figure}

% DISS-005-B1555
\subsection{Поиск осциллограмм, относящихся к конкретным терминалам/объектам}\label{sec:ch3_terminal_search}

% DISS-005-B1556
Хотя для целей создания публичного, анонимизированного датасета информация о конкретных терминалах и объектах удаляется (см. следующий раздел~3.4.3), для внутренних задач (таких как детальный анализ работы конкретного устройства, исследование специфических режимов на энергообъекте или дообучение моделей на узкоспециализированных данных) часто требуется обратная операция~-- идентификация и отбор осциллограмм, относящихся к определенным устройствам или энергообъектам.\par\vspace{0.2cm}

% DISS-005-B1557
Для этой цели были разработаны отдельные функции (find\_terminal\_hashes\_from\_json и organize\_oscillograms\_by\_terminal в SearchOscillograms), позволяющие осуществлять целевой поиск:\par\vspace{0.2cm}

% DISS-005-B1558
\textbf{Использование списка идентификаторов:} На вход подается список интересующих номеров (или других идентификаторов) терминалов.\par\vspace{0.2cm}

% DISS-005-B1559
\textbf{Сопоставление с базой данных осциллограмм:} Происходит поиск соответствий в ранее созданном JSON-файле (\_hash\_table.json), который содержит хеши осциллограмм и их исходные метаданные (имена файлов, пути). Поиск может осуществляться как по части имени файла (например, если имя содержит номер терминала), так и по пути к файлу (если структура папок отражает принадлежность к объекту/терминалу).\par\vspace{0.2cm}

% DISS-005-B1560
\textbf{Группировка осциллограмм:} Найденные осциллограммы, соответствующие заданным критериям, могут быть сгруппированы и, при необходимости, скопированы в отдельные директории для каждого терминала или объекта. Это облегчает проведение специфических исследований, например, анализ работы конкретной модели терминала РЗА или особенностей переходных процессов на определенной подстанции.\par\vspace{0.2cm}

% DISS-005-B1561
Данный функционал особенно важен при работе с большими внутренними архивами, где необходимо быстро выделить подмножество данных для решения конкретной задачи, не прибегая к повторному полному сканированию всех источников. Также это позволяет более точно задать коэффициенты нормализации для целой группы осциллограмм, провести исследование конкретного оборудования и уточнить множество других аспектов для проведения более глубоких исследований\par\vspace{0.2cm}

% DISS-005-B1562

% DISS-005-B1563
